% 普通 Markdown 仅做版面适配；公式原文保留，超宽公式的 token 副本按规则分行。
% 主模板的深度设置在此关闭，保留原稿自带的章节编号。
\setcounter{secnumdepth}{-1}
% 使用已选定的 Fandol 正文字体排数学环境中的中文，不改写公式正文。
\xeCJKsetup{CJKmath=true}
% MathJax 常用别名在此使用标准关系符实现，源包仍保留原始宏写法。
\providecommand{\lt}{<}
\providecommand{\gt}{>}
% 星号在代码字体中缺字，交由已有 CJK 正文/等宽字体显示原字符；不改代码字节。
\xeCJKDeclareCharClass{CJK}{"2605}
% Latin Modern 正文字体不含这些字符；使用现有字体显示同一 Unicode 字形。
\newfontfamily\NtwoLCircleFont[BoldFont=FandolSong-Bold.otf]{FandolSong-Regular.otf}
\newfontfamily\NtwoLProofFont{latinmodern-math.otf}
\newcommand{\NtwoLTextSymbol}[1]{%
  {\ifnum`#1="220E \NtwoLProofFont\else\NtwoLCircleFont\fi #1}%
}
% 标题中的英文词尽量整体换行，不把单词拆成跨行连字。
\pretocmd{\maketitle}{\begingroup\raggedright\hyphenpenalty=10000}{}{}
\apptocmd{\maketitle}{\endgroup}{}{}
% 小节题头至少与几行正文同页，避免长笔记中题头孤立在页底。
\pretocmd{\subsection}{\Needspace{9\baselineskip}}{}{}
\pretocmd{\subsubsection}{\Needspace{9\baselineskip}}{}{}
% 在条目开始前预留题头和正文空间；不能放入条目内部，以免推迟列表标记。
\pretocmd{\item}{\Needspace{4\baselineskip}}{}{}
\newsavebox{\NtwoLMathBox}
\newcommand{\NtwoLCurrentFormula}{}
\newcommand{\NtwoLWidthError}{%
  \typeout{N2LWIDE:\NtwoLCurrentFormula}%
  \typeout{N2LDIMS:\the\wd\NtwoLMathBox:\the\linewidth}%
  \PackageError{note2lecture}{N2L_MATH_TOO_WIDE}%
    {Formula cannot be safely wrapped within the portrait page.}%
}
\newcommand{\NtwoLDisplay}[1]{%
  \sbox{\NtwoLMathBox}{$\displaystyle #1$}%
  \ifdim\wd\NtwoLMathBox>1.12\linewidth
    \NtwoLWidthError
  \else
    \ifdim\wd\NtwoLMathBox>\linewidth
      \resizebox{\linewidth}{!}{\usebox{\NtwoLMathBox}}%
    \else
      \usebox{\NtwoLMathBox}%
    \fi
  \fi
}
% 普通独立公式保持竖版；超宽时只在明确的外层关系符处分行。
% 原公式作为完整参数保留，分行只作用于 token 副本并验证 token 不变。
\newcommand{\NtwoLMathDisplay}[2]{%
  \begingroup
  \def\NtwoLCurrentFormula{#1}%
  \sbox{\NtwoLMathBox}{$\displaystyle #2$}%
  \ifdim\wd\NtwoLMathBox>1.12\linewidth
    \NtwoLWrapDisplay{#2}%
  \else
    \[\NtwoLDisplay{#2}\]%
  \fi
  \endgroup
}
